% ============================================
% 第4章：当学生带着ChatGPT的答案走进教室
% 拱桥 · 旧系统坍塌
% ============================================

\chapter{当学生带着ChatGPT的答案走进教室}

\vspace{0.5cm}

\begin{center}
{\small\color{muted} 拱桥的秘密：把垂直的荷载，转化为水平的推力。}
{\small\color{muted} 压力不是被抵抗的，而是被利用的。}
{\small\color{muted} 石头不是被压垮的，而是被压紧的。}
\end{center}

\vspace{1cm}

\section{那一刻}

2025年10月24日，下午5点59分，桥梁风工程课。

我布置了一个课堂练习：用AI分析一段风洞数据，提取前三阶模态频率。三十秒后，一个学生举手。

\begin{shiyu-inline}
``老师，出来了。一阶0.23，二阶0.58，三阶1.04。''

我问：``这个结果，你们觉得对吗？''

教室里安静了五秒。

``呃\ldots\ldots AI 给的，应该对吧？''

\hfill ——2025年10月24日，桥梁风工程课，录音转写
\end{shiyu-inline}

那一刻，我意识到一件事：\textbf{旧系统已经坍塌了。}

旧系统是什么？旧系统是：老师拥有知识，学生需要知识，课堂是知识传递的场所。老师讲，学生记。老师考，学生答。这个系统运行了几百年，所有人都觉得理所当然。

然后ChatGPT来了。

现在，任何一个学生都可以在三十秒内得到一份完整的模态分析结果。格式规范，步骤清晰，甚至附带了解释。如果知识传递是教学的全部内容，那我已经没有存在的必要了。

\section{讲解变便宜了，什么变贵了？}

2026年5月7日，我在一个人工智能与高等教育融合的工作坊上做分享。那场分享持续了三个小时，后来被转写成了文字。中间有一段，我讲到了这个变化：

\begin{shiyu-inline}
``视频和解释随处可得，但让学习者判断'我到底懂没懂'、完成测试、调用知识解决具体问题，才是未来真正值钱的教育体验。这比'AI帮老师备课'深一层——它重新定义了课程的稀缺品。''

\hfill ——2026年5月7日，AI与高等教育融合工作坊，录音转写
\end{shiyu-inline}

这段话里有一个特别重要的词：\textbf{稀缺品}。

在旧系统里，稀缺品是\textbf{知识}。老师拥有知识，学生没有。所以老师站在讲台上，学生坐在下面。这个权力结构建立在信息的\textbf{不对称}上。

AI把这种不对称炸掉了。现在，任何人有手机就能获取任何知识。知识不再是稀缺品。\textbf{判断力才是。}

\begin{insight}
\label{insight:chap4}

\textbf{AI把``讲解''变便宜了。但把``判断''变贵了。}

过去，一个老师最值钱的能力是``讲得清楚''。未来，一个老师最值钱的能力是``设计得出让学生不得不判断的场景''。

不是``我来告诉你答案''。而是``我来制造一个你必须自己判断的情境''。
\end{insight}

\section{旧系统坍塌的三个瞬间}

旧系统不是在一天之内坍塌的。它是在三个瞬间，一点一点瓦解的。

\subsection{第一个瞬间：学生交上来的作业，AI写的}

2025年秋天，我开始在学生的作业里发现一种奇怪的``完美''。格式规范，逻辑清晰，参考文献格式正确——但所有参考文献都是编的。

我找了一个学生来问。他说：``老师，我用AI写的。但是我不知道它编了参考文献。我以为它查过。''

他不是在偷懒。他是真的不知道AI会编造参考文献。在他的认知里，AI就像一个超级搜索引擎——它给你的东西，应该是有出处的。

\textbf{旧系统的第一个裂缝：学生不再区分``知识''和``AI生成的内容''。}

\subsection{第二个瞬间：学生用AI写的作业，比老师写的教案还好}

2025年冬天，我试着用AI生成了一份涡振原理的讲义。我把它和我的教案放在一起比较。AI的版本更简洁，例子更丰富，甚至还配了流程图。

我在课堂上说了一句话，后来被转写成了文字：

\begin{shiyu-inline}
``还凑合啊，这个挺挺像那个像像那回事的啊。咱们一会儿会大概跟大家一起过一下，窝阵的这种振动现象的特点。''

\hfill ——2025年11月4日，涡振原理课，录音转写
\end{shiyu-inline}

我说``还凑合''。但心里想的是：\textbf{它比我写得好。}

\textbf{旧系统的第二个裂缝：老师发现自己的核心能力——``讲清楚''——被AI超越了。}

\subsection{第三个瞬间：学生问了一个AI回答不了的问题}

2026年春天，一个学生在课堂上问了我一个问题。不是``简支梁跨中弯矩是多少''——那个问题AI能回答。他问的是：

``老师，AI说这座桥的梁高应该取1.5米。但我用游戏试了一下，1.2米也能站住。为什么规范说1.5米？''

这个问题，AI回答不了。因为AI不知道规范里的安全系数是怎么来的。它不知道那些系数背后是几十年的工程事故和血的教训。它不知道1.5米不是``最优解''，而是``最安全的解''。

\textbf{旧系统的第三个裂缝：学生问出了一个只有人才能回答的问题。}

\section{拱桥的隐喻：压力如何变成强度}

拱桥的秘密，不是抵抗压力，而是利用压力。

你把一块石头放在拱顶上，它不往下掉——它被两边的石头挤住了。压力越大，挤得越紧。石头不是被压垮的，而是被压紧的。\textbf{压力，创造了强度。}

旧系统坍塌带来的压力，不是要摧毁教育的敌人。它是建造拱桥的石头。

\begin{itemize}
  \item \textbf{AI能生成一切→压力→教师不再靠``讲清楚''吃饭→必须设计``让学生判断''的场景。}
  \item \textbf{学生能随时获取答案→压力→课堂不再是知识传递的场所→必须变成困惑制造的场所。}
  \item \textbf{作业可以用AI完成→压力→作业不再是检验知识的手段→必须变成检验判断的手段。}
\end{itemize}

每一块压力，都对应着一块新的能力。这就是拱桥的建造方式。

\section{教师还剩什么？}

2026年7月2日，我在准备一场报告——``AI时代，教学的变与不变''。我对着录音笔说了一段话：

\begin{shiyu-inline}
``我们怎么去重新评估人的价值呢？其实我觉得比较有意思的，首先就是AI的学习方式咱们可以借鉴。它是怎么学习的呢？就是通过训练，给它大量的文本。很多时候它也不理解这东西是啥。''

\hfill ——2026年7月2日，录音转写
\end{shiyu-inline}

这段话里，我其实是在问自己：\textbf{如果AI不理解，而人理解——那``理解''到底是什么？}

后来我慢慢想明白了。理解不是``能输出正确答案''。理解是：

\begin{itemize}
  \item \textbf{能在自己错的时候，知道自己错了。}
  \item \textbf{能在一个新的情境里，判断旧的知识是否适用。}
  \item \textbf{能解释为什么选这个方案，而不是那个方案。}
  \item \textbf{能指着一样东西说：``这是我的。''}
\end{itemize}

这些，AI全做不到。不是因为它不够聪明。是因为它没有\textbf{身体}。它没有碰过混凝土，没有感受过钢筋被弯折时的阻力，没有在深夜因为一个参数调不对而摔过键盘。它没有\textbf{做过}任何事。它只是\textbf{算过}。

\textbf{教师剩下的，就是那些AI算不出来的东西。}

\begin{shiyu-note}
\label{note:chap4}

\textbf{旧系统坍塌了。不是坏事。}

旧系统本来就有问题。它假设知识的稀缺性，但知识从来都不稀缺——好奇心才是。它假设老师是知识的唯一来源，但老师从来都不是——经验才是。它假设学习是``传递''，但学习从来都不是——学习是\textbf{转化}。

AI只是把这些旧假设炸掉了。它把我们从``知识搬运工''的角色里解放了出来。现在，我们可以做真正重要的事了：\textbf{设计困惑。制造反馈。守护判断。见证骄傲。}

这些事，AI做不了。不是因为它们太难。是因为它们太\textbf{人}了。
\end{shiyu-note}

\vspace{1cm}

\begin{decision-point}
\label{decision:chap4}

\noindent\textbf{这一章，我们看到了旧系统如何坍塌，以及坍塌之后教师还剩什么。}

\vspace{0.3cm}

\noindent 下一节课，做一件事：在你讲完一个知识点之后，不要问``大家听懂了吗''。问：``AI刚才给我的答案，你们觉得对吗？为什么？''

\noindent 你会发现，学生沉默的时间变长了。然后，有人会开口。那个开口的人，说的不是``标准答案''。他说的是\textbf{判断}。
\end{decision-point}

\vspace{0.5cm}

\begin{flushright}
{\small\color{muted} 下一章：排名、积分、点赞——我们在设计``微型制度''。}
\end{flushright}